% -----------------------------------------------------------
\section{Satan, Lucifer}
% -----------------------------------------------------------
	\label{SATAN}
	\label{LUCIFER}
	
% % ----------------------
% \subsection{See also}
% % ----------------------

% % ----------------------
% \subsection{1828 American Dictionary of the English Language} % *1828
% % ----------------------
	% \label{SATAN-1828}

	% \begin{compactitem}
		% \item
	% \end{compactitem}

% % ----------------------
% \subsection{Oxford English Dictionary} % *OED
% % ----------------------
	% \label{SATAN-OED}

	% \begin{compactitem}
		% \item[]
	% \end{compactitem}

% ----------------------
\subsection{Scriptural Usage} % *SCRIPTURES
% ----------------------
	\label{SATAN-SCRIPTURES}

	% -----
	\subsubsection{Old Testament} % *OT
	% -----
		\label{SATAN-OT}
	
		% % --
		% \paragraph{KJV}
		% % --
			% \label{SATAN-OT-KJV}
	
			% \begin{tabular}{ll}
				% Total occurrences in the OT & 
			% \end{tabular}
			
			\begin{compactdesc}
				\item[{\hyperref[ISAIAH-14-12]{Isaiah 14:12}}] How art thou fallen from heaven, O Lucifer, son of the morning!  how art thou cut down to the ground, which didst weaken the nations!
					\begin{compactitem}
						\item See my notes at this verse.
					\end{compactitem}
			\end{compactdesc}
			
		% % --
		% \paragraph{Hebrew Bible}
		% % --
			% \label{SATAN-HEBREW}
		
			% %
			% \subparagraph{\texorpdfstring{\he{sample word}}{}}
			% %
				% \label{PAGA} % whatever is in step bible without periods
			
				% \begin{compactdesc}
					% \item[HALOT , PDF ] \textbf{}
						% \begin{compactitem}
							% \item[1]
						% \end{compactitem}
					% \item[BDB , PDF ] \textbf{}
						% \begin{compactitem}
							% \item[1]
						% \end{compactitem}
					% \item[DCH PDF ] \textbf{}
						% \begin{compactitem}
							% \item[1]
						% \end{compactitem}
				% \end{compactdesc}
				
				% \begin{compactdesc}
					% \item[Some OT uses] \textbf{}
						% \begin{compactdesc}
							% \item[]
						% \end{compactdesc}
				% \end{compactdesc}
			
		% % --
		% \paragraph{Septuagint}
		% % --
			% \label{SATAN-LXX}
		
			% %
			% \subparagraph{\texorpdfstring{\gr{sample word}}{}}
			% %
		
	% -----
	\subsubsection{New Testament} % *NT
	% -----
		\label{SATAN-NT}
	
		% --
		\paragraph{KJV}
		% --
			\label{SATAN-NT-KJV}
			
	
			% \begin{tabular}{ll}
				% Total occurrences in the NT & 
			% \end{tabular}
			
			\begin{compactdesc}
				\item[{\hyperref[JOHN-8-44]{John 8:44}}]
			\end{compactdesc}
			
		% % --
		% \paragraph{Greek New Testament} % *GREEK
		% % --
			% \label{SATAN-GREEK}
		
			% %
			% \subparagraph{\texorpdfstring{\gr{sample word}}{}}
			% %
				% \label{KATALLAGE}
			
				% \begin{compactdesc}
					% \item[BDAG , PDF ] \textbf{}
						% \begin{compactitem}
							% \item 
						% \end{compactitem}
					% \item[LSJ , PDF ] \textbf{}
						% \begin{compactitem}
							% \item 
						% \end{compactitem}
				% \end{compactdesc}
				
				% \begin{compactdesc}
					% \item[Some NT uses] \textbf{}
						% \begin{compactdesc}
							% \item[]
						% \end{compactdesc}
				% \end{compactdesc}
		
	% -----
	\subsubsection{Book of Mormon} % *BOM
	% -----
		\label{SATAN-BOM}
	
		% \begin{tabular}{ll}
			% Total occurrences in the BOM & 
		% \end{tabular}
		
		\begin{compactdesc}
			\item[{\hyperref[HELAMAN-6-21]{Helaman 6:21--31}}]
		\end{compactdesc}
		
	% -----
	\subsubsection{Doctrine and Covenants} % *DC
	% -----
		\label{SATAN-DC}
	
		% \begin{tabular}{ll}
			% Total occurrences in the DC & 
		% \end{tabular}
		
		\begin{compactdesc}
			\item[{\hyperref[DC1-35]{DC 1:35}}] 
			\item[{\hyperref[DC29-36]{DC 29:36--40}}] 
			\item[{\hyperref[DC76-25]{DC 76:25--49}}]
		\end{compactdesc}
		
	% -----
	\subsubsection{Pearl of Great Price} % *PGP
	% -----
		\label{SATAN-PGP}
	
		% \begin{tabular}{ll}
			% Total occurrences in the PGP & 
		% \end{tabular}
		
		\begin{compactdesc}
			\item[{\hyperref[MOSES-4-3]{Moses 4:3--21}}] Especially verses 3--4.
		\end{compactdesc}
		
% % ----------------------
% \subsection{Key Phrases} % *KEYPHRASES *PHRASES
% % ----------------------
	% \label{SATAN-PHRASES}

	% \begin{compactdesc}
		% \item[] \textbf{\#times} | list
			% % \begin{compactdesc}
				% % \item[Bible anchor verses] \phantom{.}
					% % \begin{compactdesc}
						% % \item[Romans 5:5] \gr{}
					% % \end{compactdesc}
				% % \item[Commentary] ---
			% % \end{compactdesc}
	% \end{compactdesc}
	
% % ----------------------
% \subsection{Bible Dictionaries} % *DICTIONARIES
% % ----------------------
	% \label{SATAN-BD}

% % ----------------------
% \subsection{Restoration Usage} % *RESTORATION
% % ----------------------
	% \label{SATAN-RESTORATION}

	% % -----
	% \subsubsection{Joseph Smith} % *JS
	% % -----
		% \label{SATAN-JS}
	
		% \begin{compactdesc}
			% \item[{\hyperref[KEY]{Date/Source}}]
		% \end{compactdesc}
	
	% % -----
	% \subsubsection{Brigham Young} % *BY
	% % -----
		% \label{SATAN-BY}
	
		% \begin{compactdesc}
			% \item[{\hyperref[KEY]{Date/Source}}]
		% \end{compactdesc}
	
% ----------------------
\subsection{Commentary and Studies} % *COMMENTARY *STUDIES
% ----------------------
	\label{SATAN-COMMENTARY}

	% % -----
	% \subsubsection{Key Points}
	% % -----
		% \label{SATAN-KEYS}
	
		% \begin{compactitem}
			% \item
		% \end{compactitem}
		
	% -----
	\subsubsection{Satan's power and authority}
	% -----
		\label{SATAN-power}
		
		The following passages are important:
		
		\begin{compactitem}
			\item \hyperref[MATTHEW-9-34]{Matthew 9:34}
			\item \hyperref[MATTHEW-12-26]{Matthew 12:26}---I'm thinking specifically of the ``kingdom'' word here, in connection with \hyperref[DC88-38]{DC 88:38}. See also \hyperref[LUKE-11-18]{Luke 11:17--19}.
			\item The tempation in the wilderness early in Christ's life, especially the Luke account; it's \hyperref[LUKE-4-1]{Luke 4:1--13} but verses 5--8 are particularly important.
			\item \hyperref[LUKE-12-5]{Luke 12:5}
			\item \hyperref[COLOSSIANS-1-11]{Colossians 1:11-13}
				\begin{compactitem}
					\item Verse 13 is the most important here.
				\end{compactitem}
			\item \hyperref[1NEPHI-15-35]{1 Nephi 15:35}, the earlier version, which has ``proprietor'' instead of ``preparator.'' Goes with the endowment comment below.
			\item \hyperref[2NEPHI-2-27]{2 Nephi 2:27, 29}
				\begin{compactitem}
					\item Verse 29 is super important.
					\item Read it together with 3 Nephi 28:39 below.
				\end{compactitem}
			\item See \hyperref[ALMA-34-35]{Alma 34:34--35}, the devil ``sealing you his''; see \hyperref[SEAL]{Seal}.
				\begin{compactitem}
					\item See also \hyperref[2NEPHI-9-46]{2 Nephi 9:46}.
				\end{compactitem}
			\item \hyperref[DC29-36]{DC 29:36--40}, especially verse 36; compare to \hyperref[MOSES-4-3]{Moses 4:3}
			\item \hyperref[DC76-25]{DC 76:25--31}, where it says that there was an angel of God ``who was in authority in the presence of God''; see my notes on these verses and their connection to Isaiah. See also the notes at \hyperref[DC76-31]{DC 76:31}, which allows you to connect to the oath and covenant of the priesthood at \hyperref[DC84]{DC 84}. See my notes at 26 too, for the ``man of lawlessness,'' which is extremely important.
			\item \hyperref[DC5-19]{DC 5:19}, the JSP and BC versions
			\item \hyperref[TEMPLE-ENDOWMENT-ENDOW-PRE-1990-GARDEN]{The discussion of Satan's apron} in the endowment ceremony. Satan says that the apron is ``an emblem of my power and priesthoods.'' See also \hyperref[ENDOWMENT-apron]{this note} in particular.
			\item See the bit in the \hyperref[TEMPLE-ENDOWMENT-ENDOW-PRE-1990-TELESTIAL]{endowment ceremony} about Satan being the ``god of this world.''
				\begin{compactitem}
					\item See also \hyperref[1NEPHI-15-35]{1 Nephi 15:35} above, \hyperref[DC1-35]{DC 1:35} (``the devil shall have power over his own dominion''), \hyperref[DC5-19]{DC 5:19} JSP manuscript version, 
				\end{compactitem}
			\item \hyperref[MOSES-4-1]{Moses 4:1--4} is important here too, and you can connect it to the verses in Luke 4 mentioned above.
			\item \hyperref[1NEPHI-22-26]{1 Nephi 22:26}
			\item \hyperref[2NEPHI-2-29]{2 Nephi 2:29}
			\item \hyperref[3NEPHI-6-28]{3 Nephi 6:28--30}
				\begin{compactitem}
					\item I've never made the connection with these verses before, but it says that the devil is the administrator of a certain covenant. This could be a basis for the whole way of thinking about the matter.
				\end{compactitem}
			\item \hyperref[3NEPHI-7-5]{3 Nephi 7:5} is the kind of verse, and there are others, where you can read it in the light of ideas about power and knowledge---so Satan has some sort of power stemming from some kind of knowledge and ``obedience'' to that knowledge.
				\begin{compactitem}
					\item It would help to round up other verses that talk about Satan's power this way.
					\item And I would want to connect this idea with the ones mentioned in my notes on \hyperref[GOODNESS-BOM]{the BOM on goodness}, for 3 Nephi 4:5 and other verses, where we see the effects of evil and what it is able to accomplish. I think this is one of the things that is fundamental about the difference between good and evil, or it's a way of thinking about their true nature.
				\end{compactitem}
			\item \hyperref[3NEPHI-28-39]{3 Nephi 28:39} is a very important passage, and it goes with the equally important 
			\item \hyperref[4NEPHI-1-28]{4 Nephi 1:28}
			\item \hyperref[MORMON-1-19]{Mormon 1:19}; also 2:10 here
			\item \hyperref[MORMON-5-18]{Mormon 5:18}
			\item \hyperref[ETHER-8-15]{Ether 8:15--26}
				\begin{compactitem}
					\item Also 9:5--6 here.
					\item Also 10:33.
				\end{compactitem}
			\item \hyperref[ETHER-15-19]{Ether 15:19}
		\end{compactitem}